DGFiP, IRCOM et déclarations 2042, revenus 2024 ; CCSS ; calculs de l'auteur.\par Note : différence entre la part du revenu après impôt détenue par les foyers les plus modestes avec le barème E et aujourd'hui, en pour mille, aux bornes des tranches de RFR. Positive partout : dominance de Lorenz. Elle échoue de peu entre 100 000 et 300 000 \euro{} si les revenus au PFU sont logés aussi haut que possible.
